% !TeX root = ../main.tex
\section{Структурная схема устройства}

Структурная схема разработанного тракта пространственно-временной обработки сигнала представлена на рисунке 1.

\begin{figure}[h!]
    \centering
    \begin{tikzpicture}[
        auto,
        block/.style={
            rectangle, 
            draw=black, 
            thick, 
            fill=blue!5, 
            text width=4em, 
            align=center, 
            minimum height=3em
        },
        mult/.style={
            circle, 
            draw=black, 
            thick, 
            fill=blue!5, 
            minimum size=2em, 
            inner sep=0pt
        },
        line/.style={
            draw, 
            thick, 
            -Latex
        }
    ]

    % Определение узлов (блоков схемы)
    \node[coordinate] (input) {};
    \node[block, right=1.5cm of input] (adc) {АЦП};
    \node[mult, right=2cm of adc] (mixer) {$\times$};
    \node[block, below=1.5cm of mixer] (nco) {ЦГ};
    \node[block, right=2cm of mixer] (lpf) {ФНЧ};
    \node[block, right=1.5cm of lpf] (decim) {$\downarrow D$};
    \node[block, right=1.5cm of decim] (mf) {СФ};
    \node[coordinate, right=1.5cm of mf] (output) {};

    % Соединительные линии и подписи
    \draw[line] (input) -- node[above] {$x(t)$} (adc);
    \draw[line] (adc) -- node[above] {$x_q[n]$} (mixer);
    \draw[line] (nco) -- node[right] {$s_{NCO}[n]$} (mixer);
    \draw[line] (mixer) -- node[above] {$y_{cdm}[n]$} (lpf);
    \draw[line] (lpf) -- node[above] {$y_f[n]$} (decim);
    \draw[line] (decim) -- node[above] {$y_{dec}[m]$} (mf);
    \draw[line] (mf) -- node[above] {$y_{sf}[m]$} (output);

    \end{tikzpicture}
    \caption{Структурная схема тракта обработки радиосигнала}
    \label{fig:block_diagram}
\end{figure}

На схеме обозначены следующие этапы:
\begin{itemize}
    \item \textbf{АЦП} --- аналого-цифровой преобразователь, осуществляющий субдискретизацию входного радиосигнала $x(t)$.
    \item \textbf{Смеситель ($\times$) и ЦГ} --- блок цифровой демодуляции (ЦДМ), где оцифрованный сигнал умножается на опорное колебание цифрового гетеродина для переноса на нулевую частоту.
    \item \textbf{ФНЧ} --- цифровой фильтр нижних частот, отсекающий побочные высокочастотные продукты преобразования.
    \item \textbf{$\downarrow D$} --- дециматор, понижающий частоту дискретизации.
    \item \textbf{СФ} --- блок согласованной фильтрации, вычисляющий сегментированную быструю свертку алгоритмом перекрытия со сложением.
\end{itemize}